\section{Related Work}
\label{sec:related}

\paragraph{Subliminal learning.}
\citet{cloud2025subliminal} demonstrated trait transfer through semantically
unrelated number, code, and reasoning data, and found that transfer depended on
teacher and student base-model compatibility. The phenomenon is therefore not
ordinary semantic imitation and motivates closer study of the model-specific
channels available in generated data.

\paragraph{Token-level and sparse-data accounts.}
\citet{zur2025entanglement} introduced token entanglement, measured it using
unembedding similarity and model outputs, and demonstrated subliminal prompting.
Their account supplies the prompting protocol that we scale and dissect.
\citet{schrodi2026understanding} subsequently found that global token
entanglement and logit leakage are not necessary for training-time transfer;
instead, transfer in their setting concentrates in sparse divergence tokens and
depends strongly on early trainable layers. Importantly for our novelty scope,
they already score digit-split Qwen and Gemma numbers with the autoregressive
chain rule. Our contribution is therefore not generic multi-token scoring. It is
the bidirectional first-token and within-width control that exposes a
mixed-width normalization confound. Their early-layer result concerns which
updates install a preference; our late-layer result concerns when an answer is
readable in a frozen forward pass.

\paragraph{Activation-space accounts.}
\citet{blank2026steering} interpret subliminal learning as steering-vector
distillation: a teacher system prompt induces a residual-stream direction and a
student learns an aligned direction during fine-tuning.
\citet{morgulis2026steering} show that stronger, multi-word signals can be
transmitted with trained steering vectors and report layer-localized recovery of
the teacher direction. These findings motivate looking inside the network rather
than relying only on output-token geometry. Our fixed linear readout asks when a
particular animal answer becomes readable. Our separate full-state patch tests
natural donor-to-recipient causal transmission across scale, but does not
isolate one direction or identify a training-installed signal.

\citet{wang2026data2behavior} extract layerwise mean representations from
candidate training data and causally inject them during inference to anticipate
post-training biases. A concurrent positional-bias study combines
divergence-token masking, token geometry, layerwise probing, and causal steering
across eight models \citep{positional2026}. These studies preclude a broad claim
of novelty for frozen representation analysis, layerwise subliminal probing, or
causal steering. Our narrower object is a matched 8B/70B trace of an unchanged
frozen prompting channel after its static geometry weakens.

\paragraph{Compatibility and capacity.}
In a controlled neural-network setting, \citet{brockers2026noise} show that
compatible output heads can enable subliminal transfer despite changes to hidden
layers or architecture, and derive conditions under which transfer fails. This
reinforces the need to distinguish output-boundary compatibility from the
internal computation producing an answer. Our 8B--70B comparison shows an
empirical version of this distinction: static output geometry weakens while a
late contextual readout remains similar.

\citet{madl2026channel} sharpen this distinction into a channel-conditioned
auditability map. They causally localize masked single-token transfer to
convergent vocabulary geometry in one regime, while conditional behaviors route
through model-body computation in another; their scale sweeps reach 6.9B. This
directly rules out generalizing our 70B prompting result into a universal claim
that vocabulary channels fade with scale.

The closest work therefore already covers global token entanglement,
autoregressive digit scoring, sparse divergence tokens, residual-stream
directions, causal steering, representation--behavior dissociations, and
output-head compatibility. To our knowledge from a documented search, prior
work has not combined (i) bidirectional sequence scoring with first-token and
within-width controls that diagnose the length confound, and (ii) matched
full-BF16 Llama 8B/70B observational and causal depth traces of the same frozen
prompting channel after a resolved decline in static geometry. This is a scoped
literature-search claim, not a claim of absolute priority.
